\begin{textsample}{POS Dim 2 – al – Score 28.00 – al/al\_em\_45.txt}  \label{ex:f2_pos_252}
PROFISSIONAL O Itinerário Formativo de Formação \textbf{Técnica} e Profissional fundamenta -se no Itinerário Profissional, conjunto de ocupações com identidades bem definidas no mundo do trabalho.
Pode, também, relaciona -se a determinado segmento profissional, que é a base da organização das \textbf{ofertas} para a Educação Profissional, e deve apontar para as etapas que podem ser seguidas pelos estudantes no seu processo de formação profissional, permitindo -lhe adequar, conforme suas necessidades, o Itinerário Formativo ao Itinerário Profissional, de forma a planejar sua \textbf{carreira} e projeto de vida de acordo com as demandas de um mundo do trabalho, em constante mutação.
Importante destacar algumas informações, a partir da nova Resolução CNE/CP nº 1, de 5 de \textbf{janeiro} de 2021, que Define as \textbf{Diretrizes} Curriculares Nacionais Gerais para a Educação Profissional e Tecnológica : Art. 15.
A Educação Profissional \textbf{Técnica} de Nível Médio abrange : I - \textbf{habilitação} profissional \textbf{técnica}, relacionada ao \textbf{curso} \textbf{técnico} ; II - \textbf{qualificação} profissional \textbf{técnica}, como etapa com terminalidade de \textbf{curso} \textbf{técnico} ; e III - \textbf{especialização} profissional \textbf{técnica}, na perspectiva da formação continuada.
Destaca -se ainda que os \textbf{cursos} \textbf{técnicos} devem desenvolver competências profissionais de nível tático e específico relacionados às áreas tecnológicas identificadas nos respectivos Eixos Tecnológicos.
Já os \textbf{cursos} de \textbf{qualificação} profissional como parte integrante do \textbf{itinerário} da Formação \textbf{Técnica} e Profissional do Ensino Médio será \textbf{ofertada} por meio de um ou mais \textbf{cursos} de \textbf{Qualificação} Profissional, nos termos das \textbf{Diretrizes} Curriculares Nacionais para o Ensino Médio ( DCNEM ), desde que articulados entre si, que compreendam saídas intermediárias reconhecidas pelo mercado de trabalho.
O Itinerário Formativo representa o conjunto de percursos de formação propiciados pela \textbf{instituição} de ensino, quando da \textbf{oferta} de educação profissional, dentro de cada um dos diferentes segmentos profissionais.
Deve estar organizado de forma articulada, a fim de permitir que o estudante escolha entre diferentes possibilidades de educação profissional disponíveis, que vão desde a formação inicial e continuada, de \textbf{qualificação} ( chamado de FIC ), até a \textbf{graduação}.
As \textbf{instituições} de ensino buscam nas demandas do mundo de trabalho a lógica do planejamento para a \textbf{oferta} da educação profissional.
Cada território apresenta uma \textbf{série} de processos, \textbf{instituições} e indicadores, que podem ser utilizados de fundamento para a \textbf{oferta} de \textbf{itinerários} formativos e apoiar o desenvolvimento local.
Os territórios locais, regionais e, até mesmo, os mais amplos definem exigências, relacionadas a setores produtivos e os segmentos das profissões.
Tais fatores são fundamentais para o mapeamento e \textbf{oferta} com significado do quinto \textbf{itinerário}.
Para Alagoas, no que contempla os diversos setores produtivos econômicos e sociais do território, salientamos, enquanto áreas de grande produção permanente, a Educação, Saúde e Cultura.
Ressalta -se que Alagoas, nos últimos anos, obteve expressivo investimento em hospitais nos diferentes espaços do território.
O que sugere propostas \textbf{cursos} \textbf{técnicos} nos \textbf{itinerários} formativos de ETP, voltados para os eixos tecnológicos : Ambiente e Saúde, Desenvolvimento Educacional e Social, Produção Cultural e Design.
Ainda, enquanto território, Alagoas é singular no Turismo, no Meio Ambiente e na Cultura, todavia a Economia Criativa, Cultura Empreendedora, Comunicação, tem grande destaque também, o que demonstra grande demanda para os diversos serviços vinculados.
Neste sentido, os Eixos Tecnológicos de Gestão e Negócios, Infraestrutura, Produção Alimentícia, Recursos Naturais, Turismo, Hospitalidade e Lazer, bem como Informação e Comunicação são necessidades de fomento no território de Alagoas.
Ainda citando a nova Resolução CNE/CP nº 1, no Art. 16.
Os \textbf{cursos} \textbf{técnicos} serão desenvolvidos nas formas integrada, concomitante ou subsequente ao Ensino Médio, assim caracterizadas : I - integrada, \textbf{ofertada} somente a quem já tenha concluído o Ensino Fundamental, com matrícula única na mesma \textbf{instituição}, de modo a conduzir o estudante à \textbf{habilitação} profissional \textbf{técnica} ao mesmo tempo em que conclui a última etapa da Educação Básica ; II - concomitante, \textbf{ofertada} a quem ingressa no Ensino Médio ou já o esteja \textbf{cursando}, \textbf{efetuando} -se matrículas distintas para cada \textbf{curso}, aproveitando oportunidades educacionais disponíveis, seja em unidades de ensino da mesma \textbf{instituição} ou em distintas \textbf{instituições} e redes de ensino ; III - concomitante intercomplementar, desenvolvida simultaneamente em distintas \textbf{instituições} ou redes de ensino, mas integrada no conteúdo, mediante a ação de convênio ou acordo de intercomplementaridade, para a execução de projeto pedagógico unificado ; e IV - subsequente, desenvolvida em \textbf{cursos} \textbf{destinados} exclusivamente a quem já tenha concluído o Ensino Médio.
A nova Resolução CNE/CP nº 1, que esclarece sobre as formas de \textbf{oferta} da educação profissional, também sinaliza que as propostas de \textbf{itinerários} de ETP, a partir da \textbf{instituição} de ensino, deverão observar os incisos I, II e III.
Da mesma forma observam a forma de \textbf{oferta} subsequente acontece para estudantes que já concluíram o ensino médio.
Assim é importante destacar : § 1º A \textbf{habilitação} profissional \textbf{técnica}, como uma das possibilidades de composição do \textbf{itinerário} da formação \textbf{técnico} e profissional no Ensino Médio, pode ser desenvolvida nas formas \textbf{previstas} nos incisos I, II e III deste \textbf{artigo}. § 2º Os \textbf{cursos} desenvolvidos nas formas dos incisos I e III deste \textbf{artigo}, além dos objetivos da Educação Profissional e Tecnológica, devem observar as finalidades do Ensino Médio, suas respectivas \textbf{Diretrizes} Curriculares Nacionais e outras \textbf{Diretrizes} correlatas definidas pelo Conselho Nacional de Educação, em especial os referentes à Base Nacional Comum Curricular ( BNCC ), bem como normas complementares dos respectivos sistemas de ensino. § 3º A critério dos sistemas de ensino, observadas as DCNEM, a \textbf{oferta} do \textbf{itinerário} da formação \textbf{técnica} e profissional deve considerar a inclusão de vivências práticas de trabalho, constante de \textbf{carga} \textbf{horária} específica, no setor produtivo ou em ambientes de simulação, estabelecendo parcerias e fazendo uso, quando aplicável, de instrumentos estabelecidos pela \textbf{legislação} sobre aprendizagem profissional. § 4º Na \textbf{oferta} dos \textbf{cursos} na forma dos incisos II e IV, caso o \textbf{diagnóstico} avaliativo evidencie necessidade, devem ser introduzidos conhecimentos e habilidades inerentes à Educação Básica, para complementação e atualização de estudos, garantindo, assim, o pleno desenvolvimento do perfil profissional de conclusão.
Destaca -se que a nova Resolução CNE/CP nº 1 orienta que, para \textbf{oferta} de \textbf{curso} \textbf{técnico}, em quaisquer das formas, deve ser precedida do correspondente credenciamento da unidade de ensino e de autorização do \textbf{curso} pelo Conselho Estadual de Educação de Alagoas, órgão competente do respectivo sistema de ensino.
Nas Formas de \textbf{ofertas} e composição do \textbf{itinerário} de Formação e Técnica e Profissional, importante verificar que o Aprofundamento acontece pelo \textbf{Curso} \textbf{Técnico} propriamente dito e \textbf{previsto} por \textbf{carga} \textbf{horária} estabelecida no CNCT ( de 800h, 1000h ou 1200h ) \textbf{ofertado}, ou por meio de \textbf{Curso} Técnico ( de 800h, 1000h ou 1200h ) \textbf{ofertado} através da construção de \textbf{cursos} de \textbf{Qualificação} Profissional ( FIC ) articulados entre si, mais o módulo de Formação para o Mundo do Trabalho ( articulados com os Eixos Estruturantes ), o Projeto de Vida e as \textbf{Eletivas} : Portanto, os \textbf{Itinerários} Formativos de Formação \textbf{Técnica} e Profissional, no que concerne à \textbf{carga} \textbf{horária} ( de 800h, 1000h ou 1200h ), irão variar de acordo com o \textbf{Curso} \textbf{Técnico}, que pode ser \textbf{ofertado} a partir do conjunto de FICs articuladas ou Programa de Aprendizagem Profissional.
Especificamente nestes \textbf{itinerários}, recomenda -se que os Eixos Estruturantes sejam trabalhados a partir de um Módulo de Formação para o Mundo do Trabalho, enquanto módulo ou unidades curriculares no Aprofundamento.
O objetivo do Módulo de Formação para o Mundo do Trabalho é integrar, através de atividades, projetos de integração e contextualização, os conhecimentos desenvolvidos pelas unidades curriculares do IF.
Para \textbf{itinerários} de Formação \textbf{Técnica} e Profissional, que apresentam 800 horas no \textbf{curso} \textbf{técnico}, Aprofundamento, \textbf{Eletivas} e Projetos de Vida contribuem na formação do \textbf{técnico}, portanto compõem as unidades curriculares da proposta curricular do \textbf{curso} \textbf{técnico}.
Neste modelo, a Unidade de Ensino tem maior flexibilidade, pelo menos 400 horas, para avançar na articulação com a formação geral, integração curricular e ações pedagógicas do território : Observamos que, nesta distribuição da \textbf{carga} \textbf{horária} do IF, 400 horas podem ser utilizadas para outras perspectivas e articulações, contribuindo para além da \textbf{carga} \textbf{horária} necessária nas unidades curriculares do aprofundamento, \textbf{oferta} \textbf{eletiva} e projeto de vida que concorrem para formação do \textbf{técnico} com \textbf{carga} \textbf{horária} de 800 horas.
Para \textbf{itinerários} de Formação \textbf{Técnica} e Profissional, que apresentam 1000 horas no \textbf{curso} \textbf{técnico}, Aprofundamento, \textbf{Eletivas} e Projetos de Vida contribuem na formação do \textbf{técnico}, portanto compõem as unidades curriculares da proposta curricular do \textbf{curso} \textbf{técnico}.
Neste modelo, a Unidade de Ensino tem maior flexibilidade, pelo menos 200 horas, para avançar na articulação com a formação geral, integração curricular e ações pedagógicas do território : Observamos que, nesta distribuição da \textbf{carga} \textbf{horária} do IF, 200 horas podem ser utilizadas para outras perspectivas e articulações, contribuindo para além da \textbf{carga} \textbf{horária} necessária nas unidades curriculares do aprofundamento, \textbf{oferta} \textbf{eletiva} e projeto de vida, que concorrem para formação do \textbf{técnico} com \textbf{carga} \textbf{horária} de 1000 horas.
Unidade curricular 80h - 3º ano

% matched lemmas: artigo, carga, carreira, cursar, curso, destinar, diagnóstico, diretriz, efetuar, eletivo, especialização, graduação, habilitação, horário, instituição, itinerário, janeiro, legislação, oferta, ofertar, prever, qualificação, série, técnico
\end{textsample}
